\chapter{Classical Optimization Strategies for Fragmented Aperture Antennas}
\label{app:classical}

\section{Why Classical Methods Still Matter}
\label{app:why}

This appendix is the methodology home for the optimizer-comparison
material referenced in Chapter~2.5 (\emph{How Good Is the Genetic
Algorithm?}) and Chapter~M (\emph{AI-Accelerated Fragmented Antenna
Design}). The story-level treatment in those chapters tells the
reader \emph{what we learned about classical optimizers when ground
truth was available}; this appendix tells the reader \emph{how to
reproduce, tune, or extend the methods themselves}. The split is
deliberate. Most antenna designers reaching this volume want the
qualitative results --- ``which optimizer should I reach for at my
bit count?'' --- without the bootstrap derivations or hyperparameter
sweep tables; for them, Chapter~2.5 suffices. Researchers who want
to challenge a result, port a method to a different aperture, or
understand why a published GA speedup may not be what it seems
need the methodology depth; this appendix is written for them.

Three classes of reader are anticipated. The first is the
\emph{practical antenna designer} who has $N \!\lesssim\! 30$ free
bits and a working metaheuristic and wants confirmation that her
existing pipeline is competitive --- Sec.~\ref{app:practitioner}
addresses this case directly. The second is the
\emph{research-grade benchmark consumer} who is building a method
of their own and needs an apples-to-apples baseline --- the
methodology of Secs.~\ref{app:criteria}--\ref{app:bootstrap} is
written so that any reasonable implementation can reproduce our
numbers within seed-to-seed variance. The third is the
\emph{algorithm developer} considering a new variant; here
Sec.~\ref{app:results} reports what worked and what did not at
20~bits and 24~bits, and Sec.~\ref{app:scaling} reports how those
findings transfer to 168~bits.

The methodology has matured considerably over the three exhaustive
campaigns that produced it. The earliest study \cite{Maloney1999URSI}
used an 18-bit database to compare three GA configurations; the
24-bit follow-up \cite{Maloney2013Genetic24bit} added a steady-state
variant and re-validated. The present material absorbs both earlier
studies and adds the 20-bit campaign reported in
\cite{Maloney2026Paper1}, the 24-bit campaign reported in
\cite{Maloney2026Paper2}, the 168-bit follow-up reported in
\cite{Maloney2026Paper3}, the cache-corrected accounting that has
not previously appeared in this literature, and the strict-vs-good-enough
dual-criterion framing the exhaustive ground truth uniquely enables.
The cluster-cost analysis (Sec.~\ref{app:cluster}) and the
golf-course-landscape analysis (Sec.~\ref{app:landscape}) are new
to this appendix; neither appeared in any of the three companion
papers.

\section{The Exhaustive Ground-Truth Datasets}
\label{app:datasets}

This appendix draws on three databases that have been built
specifically for the present line of work. They differ in bit count
and in whether exhaustive enumeration was possible.

\subsection{The 20-Bit Dataset (1.05~M designs)}

The 20-bit database is built on a $250 \times 250$~mm fragmented
aperture, $30 \times 30$ binary cells under two planes of symmetry
(horizontal and vertical bisectors through the feed gap), yielding
$2^{20} = 1{,}048{,}576$ free designs. Each design is simulated via
the openTDA finite-difference time-domain
solver~\cite{Yee1966FDTD,Taflove2005CompEM} at 351~frequency points
spanning 0.5--4.0~GHz, capturing the full $|S_{11}|$ and boresight-gain
spectra. The campaign was distributed across 200 AWS Graviton
\texttt{c7g.medium} spot instances and completed in a single
afternoon, producing a 4.4~GB on-disk database. The eight reference
frequency bands used throughout this book --- 0.5--0.8, 0.8--1.2,
1.2--1.6, 1.6--2.0, 2.0--2.5, 2.5--3.0, 3.0--3.5, and 3.5--4.0~GHz ---
were chosen to expose the three regimes described in
Sec.~\ref{app:results}: the easy upper bands where the
$|S_{11}| \!\le\! -10$~dB specification is met by a large fraction
of designs, the frustrated mid-bands where the basin-of-attraction
for the global optimum is small, and the balanced lower bands.

\subsection{The 24-Bit Dataset (16.8~M designs)}

The 24-bit database is built on the same physical aperture with a
finer fragment grouping under the same two planes of symmetry,
yielding $2^{24} = 16{,}777{,}216$ free designs. The frequency grid
is identical to the 20-bit dataset (351 points across 0.5--4.0~GHz)
and the eight reference bands are the same, so 20-bit and 24-bit
results are directly comparable. The 24-bit campaign was carried
out on a larger AWS fleet over roughly six elapsed hours and is
documented in \cite{Maloney2026Paper2}. At this scale the database
file is roughly 75~GB; published analyses typically work on
per-band sub-tensors to keep working-set sizes manageable.

\subsection{The 168-Bit Production Case}

At 168~bits ($25 \times 25$ quad-symmetric, 168 free bits) the
design space is $2^{168}$ and exhaustive enumeration is no longer
available. Instead the reference data takes the form of a large
multi-seed hill-climbing pool, augmented with steady-state and
memetic runs from selected reference bands. The single best fitness
discovered over the pool serves as a \emph{best-found proxy} in
place of the brute-force optimum used at 20 and 24~bits. This
substitution has consequences for both criteria: at 168~bits the
strict 0.01~dB criterion cannot be defined operationally (there is
no ``true'' fitness to be 0.01~dB away from), so we report only
Good Enough numbers measured against the best-found proxy. The
1.53$\times$ over-count factor calibrated at 20~bits
(Sec.~\ref{app:criteria}) gives a principled basis for interpreting
those numbers in context. The 168-bit dataset and the production
optimization campaign that built it are reported in
\cite{Maloney2026Paper3}.

\noindent\textit{[Placeholder, pending 168-bit follow-up paper:
detailed dataset description, HC-pool size, best-found-proxy
methodology, pool composition by method, and the cross-validation
sub-experiments that anchor the calibration.]}

\section{Two Evaluation Criteria from Exhaustive Ground Truth}
\label{app:criteria}

The exhaustive 20-bit and 24-bit databases enable two complementary
criteria for evaluating optimizer performance, each answering a
different question. The dual reporting is not optional: each
criterion is the right answer to a question the other cannot answer,
and the gap between them is itself the calibration we need to
interpret 168-bit results where only one of the two is measurable.

\subsection{Strict Basin Recovery (0.01\,dB)}

The \emph{strict} criterion declares an optimizer trial successful
if its terminal fitness lies within 0.01~dB of the brute-force
optimum on the band in question. The 0.01~dB tolerance is
deliberately tight --- below the typical FDTD-replication scatter
the database was constructed to suppress --- so that ``within 0.01~dB''
is essentially equivalent to ``in the same fitness-degeneracy
class as the global optimum.'' Note this is a fitness criterion,
not a Hamming-distance criterion: it admits fitness-degenerate
alternative designs whose bit patterns differ from the brute-force
best but whose fitness scores match within the tolerance.

The strict criterion answers the \emph{research} question, \emph{did
the algorithm recover the global basin?} It is the only criterion
that survives a hypothesis test against ``the optimizer might have
landed at a separate, lower-fitness near-optimum that the cold-start
HC also reaches by accident.'' For benchmarks intended to argue
that an optimizer reaches global optima reliably, the strict
criterion is the bar to clear.

\subsection{Good Enough Proximity (0.5\,dB)}

The \emph{Good Enough} criterion declares an optimizer trial
successful if its terminal fitness lies within 0.5~dB of the
brute-force optimum, regardless of whether the trial occupies the
brute-force basin. A design satisfying Good Enough may inhabit a
Hamming-distant near-equivalent basin --- a separate local optimum
whose fitness happens to be within 0.5~dB of the global. Good
Enough therefore answers the \emph{engineering} question, \emph{did
the algorithm deliver a design within engineering-acceptable
tolerance of the theoretical maximum?} Designs satisfying Good
Enough are deployable; designs failing it generally are not.

The 0.5~dB threshold is chosen to match published return-loss
manufacturing tolerances for the fragmented-aperture coupon: 0.5~dB
on the band-objective is, on this aperture, the difference between
``meets spec'' and ``does not.'' Other values can be used and the
analysis is straightforward to redo; the qualitative findings
reported in this appendix (HC dominates easy bands, memetic
dominates frustrated bands, pure GA never wins) are robust to
threshold choices in the 0.3--1.0~dB range.

\subsection{The 1.53$\times$ Over-Count Factor and Why Both Are Needed}

The two criteria are not redundant. An algorithm that always reaches
strict basin recovery trivially satisfies Good Enough; the converse
is generally false. At 20~bits over the 188-band procedural sweep
used in \cite{Maloney2026Paper1}, the inverse-network warm-start
reaches Good Enough on 113.2/188 bands and strict on 73.9/188 ---
a ratio of $113.2 / 73.9 \approx 1.53$. This $1.53\times$ factor
quantifies how much Good Enough \emph{over-counts} strict basin
recovery at 20~bits, and provides the bridge to the 168-bit
production case.

At 168~bits exhaustive ground truth does not exist; the strict
criterion is therefore not measurable. Only Good Enough can be
reported against a best-found proxy. The 20-bit calibration
($1.53\times$ over-count) gives a principled scaling: a Good Enough
headline at 168~bits should be interpreted as approximately
$1/1.53 \approx 65\%$ of its face value if the reader's question
is ``did this algorithm find the true optimum.'' If the reader's
question is ``did this algorithm deliver a deployable design,'' the
face-value Good Enough number is the right answer directly. Both
questions are legitimate, and the 20-bit data is what makes the
difference legible. This is the principal scientific value of
running the 20-bit benchmark in the AI-accelerated-design era:
without it, Good Enough numbers at 168~bits would not have an
anchored interpretation.

\section{Cache-Corrected Evaluation Counting}
\label{app:cache}

The published GA-vs-HC literature in pixelated antenna design has a
recurring methodological problem: the number of unique FDTD
evaluations performed by an optimizer is not the same as the number
of fitness-function calls the optimizer made. The discrepancy ---
caused by the GA's tendency to revisit previously evaluated designs
--- inflates the apparent eval count of cache-using implementations
and deflates the eval count of memory-using implementations of the
same method.

\subsection{The Seen-Set: Why Naive Eval Counts Mislead}

For an exhaustive ground-truth comparison the natural cost metric
is the number of \emph{distinct} FDTD evaluations performed by the
optimizer per trial. Naive accounting --- incrementing a counter on
every fitness call --- is distorted by the cache hit rate: in any
working implementation that wraps the FDTD solver in a
$2^{N}$-element hash table, revisiting a previously seen design
costs essentially nothing (the hash lookup returns the stored
fitness). An optimizer with high cache reuse therefore reports a
larger naive eval count than the smaller number of unique designs
it actually simulated. Unsupervised use of cache hits has been
observed in the antenna literature to inflate apparent GA speedups
when in fact the GA is benefiting from a much larger seen-set than
HC. We therefore maintain an explicit \texttt{seen} set per trial
and count only \emph{first-time} evaluations toward the optimizer's
eval budget. Cache hits cost zero in our accounting.

\subsection{Cache Hit Rates by Method}

At 20~bits the cache hit rate is approximately 15\% for HC (the
single-bit-flip neighborhood revisits a previously seen design
roughly one in seven times once a few bits have been explored).
The hit rate is substantially lower for SS-GA (typical 4--8\%
depending on recipe) because its uniform mutation tends to produce
genuinely novel offspring, and substantially higher for the
generational GA at low mutation rates (15--25\%) because crossover
on a converged population produces many near-duplicate children.

Cache hit rate alone does not determine the rank order of
methods --- but it does determine which recipes win the
hyperparameter sweep, and that ranking inversion is dramatic
enough to be worth its own subsection.

\subsection{Cache Sensitivity of Recommended Hyperparameters}

A parallel hyperparameter sweep without the seen-set cache (every
fitness call counts, including duplicates) produces \emph{the
opposite} recommended recipe in both GA families:
\begin{itemize}
\item Generational GA, cached best: $\mathrm{cxpb} = 0.2$,
      $\mathrm{mutpb} = 0.3$. Uncached best:
      $\mathrm{cxpb} = 0.3$, $\mathrm{mutpb} = 0.7$.
\item Steady-state GA, cached best: $\mathrm{cxpb} = 0.7$,
      $\mathrm{mutpb} = 0.1$. Uncached best:
      $\mathrm{cxpb} = 0.2$, $\mathrm{mutpb} = 1.0$.
\end{itemize}
The mechanism is straightforward: without a cache, the optimizer
pays full cost for every fitness call, so high-mutation recipes
that generate novel designs win. With a cache, low-mutation
recipes win because most ``new'' designs are already seen and cost
nothing. This is not a small effect: the GA at the uncached-best
recipe is more than $2\times$ slower than at the cached-best recipe
when both are run on the cached evaluator.

Reports of GA-found designs in the literature should therefore be
read in light of whether the authors maintained a seen-set cache:
the same method run with or without caching is effectively a
different algorithm, and a fair comparison requires the cached
accounting we use throughout this appendix.

\section{Bootstrap Methodology and Geometric Restart Model}
\label{app:bootstrap}

The natural metric for evaluating an optimizer pool against a
target is the number of evaluations consumed before the first trial
crosses the threshold. We report bootstrap-resampled distributions
of this quantity throughout. The methodology has three components:
the 10K-permutation bootstrap, the geometric-restart 1.4$\times$
multiplier for trials that terminate short, and the universal
confidence-scaling validated empirically across 42 band-method
combinations.

\subsection{The 10K-Permutation Bootstrap}

For each (method, band) pair we record the FDTD-equivalent
unique-eval count consumed by every trial and the best fitness it
reached. The bootstrap procedure resamples the trial pool 10{,}000
times, each time generating a uniformly random permutation of the
trials. In each ordering we accumulate per-trial eval counts until
the first trial reaches the chosen threshold (Good Enough or strict)
and record the cumulative cost. The bootstrap distribution of those
cumulative costs is the optimizer's \emph{evals-to-first-hit}
distribution; we report its median and the 5/95 percentile range.

Because every method receives the same total compute (typically
$\sim\!100{,}000$ unique-design FDTD-equivalents per band), this
normalization makes across-method comparisons fair regardless of
per-trial cost differences. HC trials cost $\sim\!60$ evals each,
so a 100K budget yields $\sim\!5{,}000$ HC trials per band; memetic
trials cost $\sim\!1{,}830$ evals each, yielding $\sim\!50$ memetic
trials per band; bootstrap places both pools on the same
evals-to-first-hit axis.

\subsection{Geometric Restart 1.4$\times$ Multiplier Validation}

Trials that terminate before crossing the threshold (a local
optimum at fitness more than 0.5~dB from the brute-force optimum,
which is most trials on the frustrated bands) are extended via
\emph{geometric restart}: the budget for the next-trial extension
is $1.4\times$ the budget consumed by the previous attempt. This
multiplier was validated empirically to give 90\% confidence at the
1~dB threshold across the 40 method-band combinations of the 20-bit
benchmark. It expresses the empirical observation that on the
frustrated bands, each successive restart of a previously failed
trial costs roughly $40\%$ more than the prior restart before the
next basin is reached.

The geometric-restart extension is required for fair accounting:
otherwise an optimizer that produces many cheap-but-unsuccessful
trials would appear to have a small evals-to-first-hit number
(because each individual trial is cheap) when in fact the cumulative
cost across the trials is large. Bootstrap with the geometric
restart adjustment treats each method's pool as a single source of
budget rather than as a sequence of independent attempts.

\subsection{Universal Confidence Scaling}

Across 42 band-method combinations spanning the 20-bit and 24-bit
campaigns, the bootstrap distributions of evals-to-first-hit are
well approximated by a geometric distribution. This gives a simple
budget rule:
\begin{itemize}
\item 90\% confidence: $3.3\times$ median evals
\item 95\% confidence: $4.3\times$ median evals
\item 99\% confidence: $6.6\times$ median evals
\end{itemize}
These multipliers can be applied to any reported bootstrap median
to translate it into a confidence-interval budget. The fit is
notably tighter for HC (which is closer to a true geometric process
because each trial is genuinely independent) and looser for memetic
(which has a strong band-dependent variance because the within-trial
HC polish is itself a partially deterministic procedure).

\subsection{When Restart-Budget Extension Is the Only Fair Eval}

For the strict criterion (0.01~dB) on the hardest bands, the 100K
per-band budget is insufficient to obtain a non-trivial
bootstrap-median: a substantial fraction of trials do not reach
strict and the geometric restart cannot extrapolate from a pool of
non-successes. In these cases the right metric is the
\emph{cumulative cold multistart} eval count --- the integer number
of evaluations consumed from cold start until the first trial in a
freshly initialized pool crosses the strict threshold. On the
hardest 20-bit bands this can reach 1--3 million evaluations.

The cumulative-cold-multistart procedure is what we use for the
cold-start reference in the warm-start speedup numbers reported in
\cite{Maloney2026Paper1}: per-band cumulative cost from cold-start
HC until the strict basin is first reached, capped at $10^7$
evaluations (a cap large enough that all eight 20-bit reference
bands resolve, the hardest being 1.6--2.0~GHz at $\sim\!3 \times
10^6$ evaluations).

\section{The Five Optimizer Families}
\label{app:methods}

We compare five optimizer families throughout. Within each family
the same hyperparameter settings are used for every band so that
comparisons reflect optimizer character rather than per-band
tuning effort. The settings themselves are the best universal
recipes from the cache-corrected sweep documented in
Sec.~\ref{app:recipes}.

\subsection{Hill Climbing (HC)}

First-improvement HC over the single-bit-flip neighborhood, started
from a uniform random bit string. On each pass, the bits are
shuffled to a random order; the algorithm flips the first bit whose
flip improves fitness, then restarts the pass from the new design.
A trial terminates when a full pass produces no improving flip
(i.e., the design is a local optimum in the $N$-bit-flip
neighborhood) and we record the best fitness and the number of
unique-design evaluations consumed. Trials are independent.

We deliberately use \emph{first-improvement} rather than
\emph{steepest-ascent} HC. First-improvement makes more passes
through the bit-shuffle order, explores more broadly, and is less
prone to early plateau lock than steepest-ascent variants on
non-convex fitness landscapes. The choice can matter by a factor
of 2--3 in evals-to-first-hit on the frustrated bands.

The single-bit-flip neighborhood is the natural choice for a binary
design space; multi-bit-flip neighborhoods (commonly used in some
combinatorial-optimization literatures) we do not benchmark here.
They would expand the neighborhood-search radius at the cost of
per-pass cost scaling with $\binom{N}{k}$ for $k$-bit flips;
preliminary results suggest the optimal radius is $k = 1$ at all
three bit counts studied.

\subsection{Generational Genetic Algorithm}

Tournament-selection GA with two-point crossover, population~32,
100 generations, tournament size~2, $\mathrm{cxpb} = 0.2$,
$\mathrm{mutpb} = 0.3$, and per-bit mutation rate
$\mathrm{indpb} = 0.5 / N_{\rm bits} = 0.025$ at 20~bits. The
hyperparameters are explained in detail in Sec.~\ref{app:recipes};
the population size $32$ is a deliberate choice (small enough to
keep cluster overhead modest, large enough to maintain diversity at
20--24 bits).

Hall-of-fame elitism keeps the best individual across generations.
This is not the DEAP default and is essential: without it, high
$\mathrm{mutpb}$ destroys the best individual every generation
(see Sec.~\ref{app:methods:deap_bug} below).

\subsection{Steady-State GA (SS-GA)}

Population~32, tournament size~3, $\mathrm{cxpb} = 0.7$,
$\mathrm{mutpb} = 0.1$, $\mathrm{indpb} = 0.025$ at 20~bits,
two-point crossover. One offspring per iteration replaces the
worst surviving individual if better. The per-trial step budget is
$\mathrm{POP} \times \mathrm{NGEN} = 3{,}200$ steady-state steps,
matched to the generational GA's per-trial evaluation budget.

SS-GA has structural advantages over the generational GA: there
are no idle nodes (one offspring per step), no straggler effect
(the next step is the next offspring, not a generation barrier),
and the population evolves gradually rather than in synchronous
generations. These advantages are documented in detail in
Sec.~\ref{app:cluster}.

Carroll's original Fortran GA driver \cite{CarrollGA} pioneered the
steady-state approach with the philosophical claim that
$\mathrm{cxpb} = \mathrm{mutpb} = 1.0$ should be the default: every
individual should get both crossover and mutation dice rolls. The
DEAP \cite{DEAP} gate convention (probabilistic gates on whether to
attempt crossover and mutation at all) is therefore extra
conservatism Carroll never used. The cached best recipe we report
above ($\mathrm{cxpb} = 0.7$, $\mathrm{mutpb} = 0.1$) sits between
the two extremes and is consistent with Carroll's spirit at the
unique-eval-cost optimum.

\subsection{SS+HC: Post-Hoc Local Polish}

SS-GA with the same settings runs to completion; the final best
individual is then locally polished by HC. The GA portion is
identical to plain SS-GA --- the HC pass is a one-time post-process
applied only to the surviving champion. This is the simplest hybrid
in the family and is interesting because the HC polish costs only
$\sim\!60$ additional evaluations and never hurts, so SS+HC is
weakly dominant over SS-GA at both 20 and 24~bits.

\subsection{Memetic: HC on Every Modified Offspring}

The memetic algorithm \cite{Moscato1989Memetic} applies HC to
\emph{every} modified offspring before tournament replacement,
throughout the run; the initial population is also HC-polished
before any GA steps execute. Mutated or crossed-over offspring
that have not been touched by any GA operator in a given step are
not re-HC-polished (they have not changed).

Memetic is the most expensive method per generation, with a
per-trial unique-eval cost of $\sim\!1{,}830$ at 20~bits (roughly
$30\times$ the HC per-trial cost). It is also the most aggressive
at exploiting basin structure: every offspring is immediately
walked to its local optimum, so the GA's population is at all
times a population of local optima, and crossover acts as a
basin-bridging operator on already-converged points. The strict
basin-recovery rates of memetic (Sec.~\ref{app:results}) reflect
this: memetic is the only one of the five families that reliably
reaches the strict 0.01~dB basin under a 5{,}000-eval per-trial
budget.

\subsection{The DEAP Elitism Bug}
\label{app:methods:deap_bug}

DEAP's HallOfFame mechanism only \emph{records} the best individual
seen across generations. It does not inject that best individual
back into the population. Combined with a non-zero $\mathrm{mutpb}$
that can mutate the current best to a worse state, the consequence
is that the best individual can be lost across a generation
transition --- a phenomenon we observed in early sweeps that
produced large negative outliers in the per-trial fitness
distribution. The fix is to explicitly replace the worst
individual with \texttt{hof[0]} after each generation's
\texttt{pop[:] = offspring} assignment. Without this fix, none of
the recipes reported in Sec.~\ref{app:recipes} produce the win
counts we report; with it, all results in this appendix are
reproducible from the published code.

\subsection{Spatial Crossover: Tested, No Significant Benefit at Small Pop}

We benchmarked a custom \texttt{cx\_portfolio} operator that
performs line, block, or disk cuts in the 2D antenna space rather
than the uniform two-point cut of the DEAP default
\texttt{cxTwoPoint}. The spatial cuts preserve adjacent-pixel
correlations through crossover, which has theoretical appeal for a
problem in which adjacent bits are physically adjacent fragments.

Empirically at population $32$, the difference between
\texttt{cxTwoPoint} and \texttt{cx\_portfolio} is not statistically
significant: 7 seeds each across 6 bands, Wilcoxon signed-rank
test, $p = 0.30$. The mechanism for the null result is that at
pop~$32$ the population converges to dominant haplotypes within
$\sim\!30$ generations and crossover mate-pairs are effectively
identical thereafter. Crossover type does not matter when both
parents are clones.

Whether spatial crossover earns its place at larger populations
(say pop~$256$ or more) is an open question; the answer almost
certainly depends on the diversity-maintenance regime and not on
the spatial cut policy per se.

\section{Hyperparameter Recipes --- How to Tune Each Method}
\label{app:recipes}

The recipes used in Sec.~\ref{app:methods} were determined by a
cache-corrected hyperparameter sweep over the same $2^{20}$
ground-truth database. We document the sweep here so the
universal-recipe choices are reproducible and the reader can judge
whether they bias the per-band ranking.

\subsection{The 20-Bit Sweep Grid (1024 Combinations $\times$ 100 Seeds)}

For each of the generational GA and the steady-state GA, we swept
four hyperparameters jointly:
\begin{itemize}
\item $\mathrm{cxpb} \in \{0.0, 0.1, 0.2, 0.3, 0.5, 0.7, 0.9, 1.0\}$
\item $\mathrm{mutpb} \in \{0.0, 0.1, 0.2, 0.3, 0.5, 0.7, 0.9, 1.0\}$
\item $\mathrm{indpb} \in \{0.5, 1, 1.5, 2, 2.5, 3, 4, 6\}$ bits / $N_{\rm bits}$
\item $\mathrm{tournsize} \in \{2, 3\}$
\end{itemize}
This produces $8 \times 8 \times 8 \times 2 = 1{,}024$ combinations
per method. Each combination was evaluated with 100~random seeds to
a 0.5~dB Good Enough target on each of the eight reference bands,
with the same cache-corrected unique-eval accounting used in
Sec.~\ref{app:methods}. Total: $1{,}024 \times 100 \times 8 \approx
8.2 \times 10^5$ runs per method.

\subsection{Best Universal Recipes}

We rank recipes by the geometric mean of the bootstrap-median
unique-eval count across all eight bands, normalized to HC's
per-band median. The geometric mean rewards recipes that are
competent everywhere over recipes that are excellent on one band
and poor on another.

\begin{table}[h]
\centering
\caption{Best universal recipes from the cache-corrected
hyperparameter sweep, 20-bit fragmented aperture. The ratio is the
geometric mean of (method evals)/(HC evals) across all eight
reference bands.}
\label{tab:best_recipes}
\small
\begin{tabular}{lcccccc}
\hline
Method & cxpb & mutpb & indpb & tourn & pop & geo\,mean / HC \\
\hline
GA gen.   & 0.2 & 0.3 & $0.5/N$ & 2 & 32 & 0.91$\times$ \\
SS-GA     & 0.7 & 0.1 & $0.5/N$ & 3 & 32 & 0.87$\times$ \\
\hline
\end{tabular}
\end{table}

The two recipes are stark in their divergence. The generational GA
prefers \emph{low} crossover and \emph{moderate} mutation; the
steady-state GA prefers the opposite, \emph{high} crossover and
\emph{low} mutation. The asymmetry has a mechanistic explanation
already given in Sec.~\ref{app:cache} above: with
the cache active, low-mutation steady-state recipes win because
their offspring are mostly novel (the population turns over slowly
enough that crossover acts on diverse parents), while generational
GAs prefer modest mutation to keep generations producing distinct
candidates.

\subsection{Tournament Size and Mutation Rate Observations}

Two further observations from the full sweep are worth recording:
\begin{itemize}
\item Tournament size~3 dominates tournament size~2 for the SS-GA
      (reverse for the generational GA). The smaller pop in SS-GA
      benefits from increased selection pressure; the generational
      GA already has more selection pressure than it can usefully
      apply.
\item Setting $\mathrm{mutpb} = 0$ is catastrophic for both methods
      regardless of crossover --- the GA can only explore via
      mutation; pure crossover restricts the search to subspaces
      spanned by the initial population, which is not enough to
      escape early lock.
\item Tournament size~3 on pop~$32$ crushes population diversity
      within $\sim\!30$ generations. By that point most offspring
      are clones of the current best; crossover then mates
      identical parents and produces no new information. Larger
      populations defer this collapse proportionally.
\end{itemize}

\subsection{The 168-Bit Recipe: Why 20-Bit Recipes Don't Transfer}

At 168~bits, the 20-bit cached recipe ($\mathrm{indpb} = 0.5/N = 0.025$)
remains the per-bit mutation rate. Naive transfer of the 20-bit
\emph{recipe} (the absolute $\mathrm{indpb}$ value $0.025$ ported
verbatim) is catastrophic: it flips $0.025 \times 168 \approx 4.2$
bits per mutation, well above the basin-bridging optimum for the
168-bit landscape. A 2D sweep over $\mathrm{mutpb}$ and
$\mathrm{indpb}$ at 168~bits identifies the best combination as
$\mathrm{mutpb} = 0.40$ with $\mathrm{indpb}$ scaled to flip
$\sim\!2$ bits per mutation event (i.e., $\mathrm{indpb} = 2/168
\approx 0.012$). $\mathrm{mutpb} = 0.30$ is uniformly mediocre at
168~bits and $\mathrm{mutpb} = 0.40$--$0.50$ is the practical
window.

The general lesson: at higher bit counts the optimal
\emph{per-mutation-event bit count} stays small (1--2 bits) while
the \emph{per-individual mutation probability} stays moderate
(30--50\%); the two cannot be transferred independently. Recipes
quoted as $(cxpb, mutpb, indpb)$ triples without an explicit bit
count are ambiguous.

\subsection{Recommended Starting Recipes Per Bit Count}

\begin{table}[h]
\centering
\caption{Recommended starting recipes for the three bit counts
studied. ``Bits/mut'' is the expected number of bits flipped per
mutation event when $\mathrm{mutpb}$ fires. Tournament size 2
unless otherwise noted.}
\label{tab:recipes_summary}
\small
\begin{tabular}{lcccccc}
\hline
Bits & Method & cxpb & mutpb & indpb & bits/mut & pop \\
\hline
20  & GA gen. & 0.2 & 0.3 & $0.5/N$    & 0.5 & 32 \\
20  & SS-GA   & 0.7 & 0.1 & $0.5/N$    & 0.5 & 32 (t=3) \\
24  & GA gen. & 0.2 & 0.3 & $0.5/N$    & 0.5 & 32 \\
24  & SS-GA   & 0.7 & 0.1 & $0.5/N$    & 0.5 & 32 (t=3) \\
168 & GA gen. & 0.5 & 0.4 & $2.0/N$    & 2   & 32 \\
168 & SS-GA   & 0.7 & 0.4 & $2.0/N$    & 2   & 32 (t=3) \\
\hline
\end{tabular}
\end{table}

The recipes for memetic and SS+HC are inherited from SS-GA (both
hybrids wrap the SS-GA core); HC has no recipe to tune. The 168-bit
entries are populated from the 2D sweep referenced in
\cite{Maloney2026Paper3}.

\section{Per-Band Optimizer Comparison Results}
\label{app:results}

The per-band optimizer comparison is the central empirical result
of this appendix. We report it at both criteria with full per-band
tables; the compressed versions appearing in
\cite{Maloney2026Paper1}~\S5 and \cite{Maloney2026Paper2}~\S3 are
extracted from the data here.

\subsection{The Master Comparison at 20~Bits --- Good Enough and Strict}

Under the Good Enough criterion (0.5~dB target), the win-count
summary across the eight reference bands is:
\begin{itemize}
\item HC: 3/8 (the easy bands 2.5--3.0, 3.0--3.5, 3.5--4.0~GHz)
\item SS-GA: 2/8 (the balanced lower band and one mid-band crossover)
\item Memetic: 2/8 (the frustrated mid-bands 1.2--1.6, 1.6--2.0~GHz)
\item SS+HC: 1/8 (a narrow margin)
\item Pure GA: 0/8
\end{itemize}
Under the strict criterion (0.01~dB target), the win-count summary
inverts dramatically:
\begin{itemize}
\item Memetic: 5.85/8 (a fractional count reflecting partial wins on
      shared bands; effectively dominant on every band)
\item HC: 0.18/8
\item SS-GA: 0.31/8
\item SS+HC: 0.33/8
\item Pure GA: 0/8
\end{itemize}
Memetic is the only method reaching strict basin recovery on more
than 5\% of trials on 7 of 8 bands and more than 87\% on 6 of 8;
HC, SS-GA, and SS+HC reach strict on $\leq\!5\%$ on 7 of 8 bands.

The full per-band reach-rate table is reproduced in
Table~\ref{tab:strict_reach_20bit}, and the two-criteria visual
summary in Fig.~\ref{fig:master_comparison_20bit}. Panel~A is the
bootstrap-median evals-to-Good-Enough by method and band (note the
log y-axis); Panel~B is the per-band strict reach rate.

\begin{figure}[h]
\centering
\includegraphics[width=\linewidth]{/Users/jim.maloney/github-clones/fragmented-aperture-book/images/appendixB/master_comparison_20bit_book.pdf}
\caption{Per-band optimizer comparison across the $2^{20}$ database
under both criteria. \textbf{Panel A (Good Enough, 0.5\,dB):}
bootstrap-median FDTD-equivalent evaluations to reach within
0.5\,dB of the brute-force optimum. Memetic costs $\sim\!30\times$
HC per trial, visible as a roughly flat $\sim\!1{,}800$-eval bar
across bands; HC, SS-GA, and SS+HC vary by band from $\sim\!200$ to
$\sim\!5{,}000$ evals. \textbf{Panel B (Strict, 0.01\,dB):} per-band
reach rate (\% of trials) recovering the brute-force basin within
the same $\sim\!5{,}000$-eval per-trial budget. HC, SS-GA, and
SS+HC reach strict on $\leq\!5\%$ on 7 of 8 bands; memetic reaches
strict on $\geq\!46\%$ on 7 of 8 bands and $\geq\!87\%$ on 6 of 8.}
\label{fig:master_comparison_20bit}
\end{figure}

\begin{table}[h]
\centering
\caption{Per-band reach rate (\% of trials) at strict (0.01\,dB)
vs Good Enough (0.5\,dB) thresholds, across HC, SS-GA, SS+HC, and
memetic at the 5{,}000-eval per-trial budget. 20-bit data, HC pool
of 5{,}000 trials/band, others 500 trials/band. The mean row is
the per-method average across bands.}
\label{tab:strict_reach_20bit}
\small
\begin{tabular}{lrrrrrrrr}
\hline
       & \multicolumn{2}{c}{HC} & \multicolumn{2}{c}{SS} & \multicolumn{2}{c}{SS+HC} & \multicolumn{2}{c}{Memetic} \\
Band (GHz) & strict & good & strict & good & strict & good & strict & good \\
\hline
0.5--0.8   & 0.3  &  5.8 & 2.0  & 14.6 & 2.0  & 15.2 & \textbf{93.6} &  95.8 \\
0.8--1.2   & 0.1  &  6.0 & 0.2  & 22.4 & 0.2  & 23.8 & \textbf{34.2} & 100.0 \\
1.2--1.6   & 0.5  &  1.3 & 0.4  &  1.2 & 0.4  &  1.2 & \textbf{95.0} &  95.0 \\
1.6--2.0   & 1.1  &  1.3 & 1.2  &  1.8 & 1.2  &  1.8 & \textbf{62.4} &  62.4 \\
2.0--2.5   & 1.3  &  4.7 & 4.8  & 13.6 & 5.6  & 15.0 & \textbf{98.6} & 100.0 \\
2.5--3.0   & 3.3  &  4.5 & 3.8  &  5.2 & 3.8  &  5.2 & \textbf{87.2} &  99.4 \\
3.0--3.5   & 0.8  & 17.1 & 1.6  & 26.2 & 1.8  & 31.6 & \textbf{46.8} & 100.0 \\
3.5--4.0   & 10.2 & 19.3 & 17.6 & 31.6 & 18.0 & 35.2 & \textbf{100.0} & 100.0 \\
\hline
mean       & 2.2  & 7.5  & 4.0  & 14.6 & 4.1  & 16.1 & \textbf{77.2} & 94.1 \\
\hline
\end{tabular}
\end{table}

\subsection{The Master Comparison at 24~Bits --- Regime Confirmation}

The 24-bit campaign~\cite{Maloney2026Paper2} confirms the regime
structure observed at 20~bits. Re-rendered at book scale in
Fig.~\ref{fig:master_comparison_24bit}, the 24-bit master
comparison shows that Memetic remains the only method reliably
reaching the strict 0.01\,dB basin, with per-band reach $\geq\!54\%$
on 6 of 8 bands and 100\% on the two extreme bands (0.5--0.8 and
3.5--4.0\,GHz). The single qualitative difference from 20~bits is
that the 0.8--1.2\,GHz band becomes uniformly hard at 24~bits ---
all five methods reach strict on $<\!1\%$ of trials --- a feature
visible in Panel~B and consistent with the slightly tighter
brute-force optimum on that band ($-0.20$\,dB vs $-0.67$\,dB at
20~bits, a 0.47\,dB shift). Under Good Enough the 24-bit
bootstrap-medians span a slightly wider dynamic range than at
20~bits (Panel~A), but the rank order is preserved on every band.

\begin{figure}[h]
\centering
\includegraphics[width=\linewidth]{/Users/jim.maloney/github-clones/fragmented-aperture-book/images/appendixB/master_comparison_24bit_book.pdf}
\caption{Per-band optimizer comparison across the $2^{24}$ database,
same layout as Fig.~\ref{fig:master_comparison_20bit}.
\textbf{Panel A (Good Enough, 0.5\,dB):} bootstrap-median
FDTD-equivalent evaluations to reach within 0.5\,dB. The
1.2--1.6 and 1.6--2.0\,GHz mid-bands at 20~bits relocate to a
single 0.8--1.2 GHz extreme at 24~bits where all four classical
methods need $\gtrsim\!5{,}000$ evals. \textbf{Panel B (Strict,
0.01\,dB basin recovery):} per-band reach rate (\% of trials).
Memetic dominates strict on 7 of 8 bands ($\geq\!54\%$); HC, SS-GA,
and SS+HC reach $\leq\!2\%$ on the same 7 bands.}
\label{fig:master_comparison_24bit}
\end{figure}

\subsection{The Three Regimes: Easy, Frustrated, Balanced}

Three mechanistically distinct regimes appear consistently across
20-bit and 24-bit data:

\textbf{Optimizer-easy bands (2.5--4.0~GHz).} HC dominates by
$1.5\times$ to $3\times$ over the next-best method. At these
frequencies the $|S_{11}| \le -10$~dB specification is met by a
large fraction of the design space; basins are large and shallow,
and HC's commitment to climb every uphill neighbor wins
immediately. GA, SS-GA, and SS+HC pay tournament and crossover
overhead they cannot recoup.

\textbf{Frustrated mid-bands (1.2--2.0~GHz).} Memetic dominates by
$1.6\times$ to $2\times$ over HC. The brute-force optima at these
bands are only marginally below the $-10$~dB $|S_{11}|$ line, and
the landscape is a ``golf course'' (Sec.~\ref{app:landscape}):
basin enumeration finds 1{,}000 to 2{,}900 distinct local optima
across the band's full landscape, of which only a small subset
($\sim\!1\%$) lie within 1~dB of the global optimum. HC's pool of
5{,}000 random starts saturates --- only 1--3\% of starts reach the
target basin --- while memetic's habit of climbing every new GA
offspring extracts work from every individual the GA samples.

\textbf{Balanced bands (0.5--1.2, 2.0--2.5~GHz).} SS-GA and SS+HC
win by narrow margins ($\leq\!10\%$ over HC). The basins are
intermediate in size and number; the steady-state replacement
policy avoids the generational GA's wasteful re-evaluation of
pop-wide near-duplicates and matches HC's local efficiency once
the GA has located a basin.

The regime structure is robust to bit count: each regime exists at
both 20 and 24~bits with the same per-band assignments. We expect
it to persist at 168~bits, but the prediction is testable only
indirectly because exhaustive ground truth is unavailable at that
scale.

\subsection{Win-Count Summaries --- Good Enough and Strict}

The Good Enough criterion's win-count summary HC 3/8, SS-GA 2/8,
Memetic 2/8, SS+HC 1/8, GA 0/8 (at 20~bits) and HC 4/8, SS-GA 2/8,
Memetic 1/8, SS+HC 1/8, GA 0/8 (at 24~bits) is the result most
practitioners want. Under Good Enough, no single method dominates,
the per-band winner depends on the regime, and HC is a reasonable
default.

The strict criterion's win-count summary Memetic 5.85/8 (others
$\leq\!0.33/8$) at 20~bits is qualitatively different. Under
strict, memetic is the only reliable method --- the cost on its
per-trial budget is $\sim\!30\times$ HC's, but the cost per
\emph{recovered basin} is $313$ evals (vs HC's $333$), so memetic
is within $5\%$ of HC on the strict-cost-per-basin metric. The
practical distinction is that memetic converts trials into wins
while HC, SS-GA, and SS+HC convert trials into ``close-but-not-in-basin''
designs.

\subsection{Per-Band Reach-Rate Tables --- Full Strict + Good Enough}

The 20-bit per-band reach-rate table is reproduced as
Table~\ref{tab:strict_reach_20bit} above. The 24-bit version is
qualitatively identical; the principal numerical difference is
that the dynamic range across bands narrows slightly (the easy
bands are not quite as easy at 24~bits, the frustrated bands are
slightly less frustrated). Memetic's dominance under strict
remains at all eight bands.

The book-scale 20-bit master comparison figure is
Fig.~\ref{fig:master_comparison_20bit} (above, in \S~B.8.1). The
24-bit sibling figure (which renders qualitatively the same regime
structure) is deferred pending re-render of the 24-bit
\texttt{.npz} data from \cite{Maloney2026Paper2} at book column
width.

\section{Cluster-Cost Analysis}
\label{app:cluster}

The single-machine bootstrap-eval-count metric of
Sec.~\ref{app:bootstrap} understates the cost of the generational
GA on a real cluster. On a one-node-per-population-member
deployment --- the natural cluster shape for an embarrassingly
parallel GA --- the generational GA pays three additional costs
that do not appear in the cached eval count.

\subsection{The Three Strikes Against Generational GA on Clusters}

\textbf{Strike one: idle nodes.} Within a generation, individuals
that have not been touched by any GA operator (no crossover, no
mutation) do not need new FDTD simulation; they are unchanged from
the previous generation and their fitness is cached. On a cluster,
those nodes sit idle while the others compute, but they still
incur per-hour spot-instance cost. The fraction of idle nodes per
generation is determined by the recipe: high-mutation,
high-crossover recipes touch most individuals and have small idle
fractions, low-rate recipes touch few and have large idle
fractions.

\textbf{Strike two: the straggler effect.} A generation barrier
synchronizes the population: the next generation cannot begin until
the slowest FDTD evaluation in the batch completes. On the
fragmented aperture, FDTD evaluation cost varies by a factor of
$\sim\!1.8$ across bands (high-Q designs at 0.5--0.8~GHz take
$\sim\!6$~s per evaluation; broadband designs at 1.2--1.6~GHz take
$\sim\!3.3$~s) due to early-termination cutoff differences.
Generational GA pays the maximum per generation; HC and SS-GA pay
the average per step. On a 32-node cluster this is a $\sim\!1.8\times$
multiplier on the generational GA's wall time, with corresponding
spot-cost inflation.

\textbf{Strike three: FDTD cost is design-dependent.} Beyond the
per-band variation, individual designs within a band can vary in
FDTD cost by a similar factor depending on early-termination
behavior. Generational GA pays the maximum within each batch; the
straggler effect compounds with the cross-band variance.

\subsection{Idle-Fraction by Recipe}

At 24~bits across the eight reference bands, the per-recipe idle
fractions are:
\begin{itemize}
\item Carroll-style ($\mathrm{cxpb} = \mathrm{mutpb} = 1.0$): 0\%.
      Every individual is touched, no idle nodes.
\item DEAP cached-best for generational ($\mathrm{cxpb} = 0.2$,
      $\mathrm{mutpb} = 0.3$): 21\%. About one in five nodes is
      idle each generation.
\item Conservative ($\mathrm{cxpb} = 0.1$, $\mathrm{mutpb} = 0.1$):
      89\%. The recipe minimizes seen-set churn but pays for it
      in cluster idle.
\end{itemize}
The Carroll-style recipe was empirically derived for a different
optimization objective (eval-cost minimization on a single
machine), but it has the side benefit of zero cluster idle. This
is a strong argument for the Carroll philosophy on clusters
specifically.

\subsection{The Straggler Effect: 1.8$\times$ FDTD Speed Spread by Band}

The factor-1.8 cross-band FDTD speed spread is roughly stable across
the 20-bit and 24-bit campaigns. The mechanism is the early-termination
criterion of the FDTD solver: a design with a long ring-down (typically
high-Q, low-band designs) needs more time steps before the time-domain
signal has decayed below the early-termination threshold and the
simulation can stop. Broadband mid-bands ring down quickly and
terminate early.

\subsection{Why Steady-State and HC Avoid These Problems}

Steady-state GA produces one offspring per step; there is no
generation barrier and therefore no straggler effect, and the
single offspring is always touched by at least one operator so
there are no idle nodes. The cluster shape is one node per
``current offspring'' rather than one node per population member;
the steady-state GA at pop~$32$ on a 32-node cluster uses 31 nodes
to store the population (effectively free) and 1 node to compute
the next offspring (busy).

HC is sequential by nature and runs on a single node. The
embarrassingly parallel deployment is one node per HC trial, not
one node per HC step; the per-trial cost is independent of cluster
size and the trial-pool is fully parallelizable across as many
nodes as one has.

The cluster reality is therefore that the generational GA is
\emph{disqualified} for production FDTD-driven optimization at the
scale of practical antenna design. Steady-state, HC, and memetic
hybrids are the only viable choices when each fitness evaluation
takes seconds to minutes.

\section{The Fitness Landscape Is a Golf Course}
\label{app:landscape}

The per-band rankings of Sec.~\ref{app:results} have a single
underlying cause: the fragmented-aperture fitness function as a
function of the binary mask is best modeled not as a smoothly
multimodal surface but as a ``golf course'' --- a small number of
large basins surrounded by a sea of shallow local optima.

\subsection{Steepest-Ascent Basin Mapping}

The 24-bit database makes basin enumeration tractable. For each
band we ran steepest-ascent HC from every one of the
$2^{24} = 16{,}777{,}216$ designs and recorded the local optimum
each starting point reached. The endpoints partition the design
space into basins; the basin count and basin size distribution per
band are the principal landscape diagnostics.

Across the eight bands at 24~bits we find between 1{,}000 and
2{,}900 distinct local optima per band. The largest basin holds only
2--5\% of designs per band; there is no dominant funnel. The basin
entropy is 8--10 bits (out of a maximum of 24, achieved if every
design were its own basin), indicating moderate structure: good
designs share roughly 14 of 24 bits, but past that point the
landscape is highly multimodal.

Fig.~\ref{fig:landscape_summary_24bit} summarizes the basin
statistics per band: basin count, largest-basin fraction, target
basin coverage at 0.5\,dB, and the geometric-mean restart cost
inferred from coverage. The corresponding 20-bit version is
qualitatively similar but with fewer basins (300--900 per band) and
larger target coverage (3--40\% per band).

\begin{figure}[h]
\centering
\includegraphics[width=\linewidth]{/Users/jim.maloney/github-clones/fragmented-aperture-book/images/appendixB/landscape_summary_24bit.pdf}
\caption{Basin-mapping summary across the eight 24-bit reference
bands. Each panel groups per-band statistics from steepest-ascent
HC over all $2^{24}$ designs: number of distinct local optima,
largest-basin mass (\% of designs), basin-coverage of the
0.5\,dB target region, and predicted geometric-mean restart cost.
The mid-bands 1.2--2.0\,GHz exhibit the lowest target coverage
($\sim\!1\%$), explaining the high Good-Enough costs in
Fig.~\ref{fig:master_comparison_24bit}.}
\label{fig:landscape_summary_24bit}
\end{figure}

\subsection{The 90\% Elite Paradox}

The high basin count is not the obstacle. Approximately 90\% of
random starts reach a top-1\% peak (within the best $\sim\!170{,}000$
designs by fitness across the database). The obstacle is the
restricted reach: only 1--34\% of starts reach within 0.5~dB of the
global best, depending on band. On the frustrated mid-bands the
target-basin coverage is only 1--2\%, which is the lottery-ticket
regime that makes HC restart-inefficient and memetic worthwhile.

The paradox is that ``almost every start reaches a peak'' and
``almost no start reaches the global'' coexist on the same
landscape. The resolution is that there are many high-quality
peaks (the top 1\% covers 90\% of basins by mass) but few
\emph{globally best} peaks (the within-0.5-dB region covers
1--34\% by mass, varying band by band).

\subsection{Crossover Benefit Ratio}

Crossover is sometimes argued to be destructive on fragmented-aperture
landscapes (two near-optimum parents in different basins produce a
child in neither basin). We measured this directly: across the 24-bit
landscape, what fraction of offspring produced by two-point
crossover beats both parents on fitness? The answer is 15--47\%
depending on parent quality and band, with a typical value around
25\%. Even among top 0.1\% designs (where one might expect
diminishing returns), the crossover benefit ratio remains 15--37\%,
well above the ``destructive'' threshold of 15\% (the rate at which
random bit-flips of a near-optimum would beat that near-optimum).

Crossover is therefore constructive, not destructive. The reason
pure GA still loses is different: the GA pays $\sim\!32$ evaluations
per generation for $\sim\!6$--$7$ crossover improvements (about
$20\%$ of pop~$32$), while HC pays $24$ evaluations per pass for
$\sim\!3$--$8$ steepest-direction improvements. Per evaluation
spent, HC's improvements are larger and more reliable; the GA's
population overhead exceeds the crossover benefit per evaluation.

\begin{figure}[h]
\centering
\includegraphics[width=\linewidth]{/Users/jim.maloney/github-clones/fragmented-aperture-book/images/appendixB/crossover_benefit_24bit.pdf}
\caption{Crossover-benefit ratio across the 24-bit landscape, by
parent-quality stratum and band. ``Beats both parents'' = fraction
of two-point crossover offspring with fitness exceeding both
parents on the band's exhaustive-database fitness. The 15\%
``destructive'' floor (random-flip rate) is exceeded on every band
and at every parent quality, including the top 0.1\% stratum.}
\label{fig:crossover_benefit_24bit}
\end{figure}

\subsection{HC vs SS Efficiency Decomposition}

Learning curves at fixed evaluation budgets (25, 50, 100, 200, 500,
1000, 2000, 5000 evals) show that HC and SS solve different
problems on the golf-course landscape:
\begin{itemize}
\item At 25--50 evals, HC leads by 1.0--1.7 dB across the eight
      24-bit bands. HC has completed 1--2 full passes and found a
      local peak; SS has barely initialized its population.
\item At $\sim\!100$ evals, the curves cross. SS catches up as
      its population converges and crossover benefit kicks in;
      HC is already plateaued at its local peak.
\item At 200+ evals, SS leads by 0.3--1.6 dB. SS continues
      improving via crossover and population diversity; HC is
      flat (it converged at $\sim\!50$--$100$ evals).
\end{itemize}
This pattern is consistent on all 8 bands at 24~bits and on all 6
bands at 20~bits. The cross-over point is universal: $\sim\!75$--$100$
evals across the studied bit counts, slightly earlier at 20~bits
because HC passes are shorter.

The decomposition explains why HC wins the master comparison
nevertheless. HC solves ``find a peak, fast,'' which is the right
problem when restarts are cheap (random initial designs cost nothing
to generate, and the per-trial cost is small). SS solves ``find a
\emph{good} peak,'' which is the right problem when restarts are
expensive. At 24~bits with $\sim\!5$~s FDTD per evaluation,
restarts are cheap; HC wins. At 168~bits with $\sim\!15$~s FDTD per
evaluation, restarts are expensive enough that SS's per-trial
quality advantage may outweigh its startup cost.

\begin{figure}[h]
\centering
\includegraphics[width=\linewidth]{/Users/jim.maloney/github-clones/fragmented-aperture-book/images/appendixB/efficiency_decomposition_24bit.pdf}
\caption{Learning curves at fixed evaluation budgets (25, 50, 100,
200, 500, 1{,}000, 2{,}000, 5{,}000 evals) for HC and steady-state
GA across the 24-bit landscape. HC leads at 25--50 evals on every
band; the curves cross near 75--100 evals; SS leads at 200+ evals.
The same pattern holds across all 8 bands at 24~bits and on the 6
20-bit bands.}
\label{fig:efficiency_decomposition_24bit}
\end{figure}

\subsection{Method-Selection Rule by Basin Coverage}

The full per-band, per-method analysis suggests a simple selection
rule based on the basin coverage of the target threshold:
\begin{itemize}
\item Less than 3\% target coverage $\to$ Memetic (the only method
      that reliably finds the rare basin).
\item 3--10\% coverage $\to$ HC or SS+HC (cheap restarts work; the
      HC polish at the end of SS+HC adds insurance at small cost).
\item Greater than 10\% coverage $\to$ SS (population diversity
      helps cheaply; the per-trial quality advantage matters when
      the target is a wide basin).
\end{itemize}

\begin{figure}[h]
\centering
\includegraphics[width=\linewidth]{/Users/jim.maloney/github-clones/fragmented-aperture-book/images/appendixB/basin_coverage_vs_threshold.pdf}
\caption{Basin coverage vs target-fitness threshold across the
eight 20-bit reference bands. Each band's curve gives the fraction
of designs with fitness within the threshold of the brute-force
optimum --- the ``size'' of the target basin under the chosen
criterion. At the 0.5\,dB Good Enough threshold (vertical
reference), coverage spans roughly two decades (1--34\%), partitioning
the bands into the three selection-rule regimes.}
\label{fig:basin_coverage}
\end{figure}
The rule is validated at both 20 and 24~bits with no exception. It
should transfer to higher bit counts to the extent that the
golf-course landscape character persists; preliminary 168-bit data
suggests it does (see Sec.~\ref{app:scaling}).

\section{Scaling 20 $\to$ 24 $\to$ 168 Bits}
\label{app:scaling}

The natural objection to a 20-bit benchmark is that practical
fragmented apertures have many more free bits. We address the
transfer of the 20-bit and 24-bit findings to 168~bits directly
here.

\subsection{What Stays the Same}

The regime structure (easy / frustrated / balanced) persists at all
three bit counts. The two-criteria utility is preserved (Good Enough
and strict answer different questions, and the strict-vs-Good-Enough
gap calibrates Good Enough numbers at production scale). The
cache-correction mechanism is essential at all three bit counts
because the cache hit rate is recipe-determined and recipe-dependent
in the same way at every scale.

\subsection{What Changes}

HC's neighborhood size grows linearly with bit count: at 168~bits a
single pass is 168~evaluations rather than 24 or 20. This makes
random restarts more expensive in absolute terms. Steady-state and
memetic startup costs also grow (population initialization at 168
bits costs $32 \times 168 = 5{,}376$ evaluations, vs $32 \times 24 =
768$ at 24~bits). The optimal $\mathrm{indpb}$ stays at $\sim\!0.5$
bits per mutation event at 20~bits but rises to $\sim\!2$ bits per
mutation event at 168~bits (Sec.~\ref{app:recipes}).

\subsection{The 168-Bit Surprise: The Landscape Is EASIER, Not Harder}

A Phase~1 HC168 campaign~\cite{Maloney2026Paper3} ran 36 random-restart
HC trials across six 168-bit reference bands
(0.5--0.7, 0.7--1.0, 1.0--1.4, 1.4--1.9, 1.9--2.4, and 2.4--3.0\,GHz);
six independent seeds per band, distributed across two initialization
strategies (25\% and 70\% metal-fill initial bit strings). The
per-band median trial cost is 520--1650 evaluations to local
convergence, comparable to the 20-bit HC pool's
$\sim\!60$--$300$~evals-per-trial scaled by the larger
neighborhood size (168 vs 20). The expected catastrophic scaling
under a curse-of-dimensionality reading --- which would predict
per-trial costs growing exponentially in bit count --- is absent.

\begin{figure}[h]
\centering
\includegraphics[width=\linewidth]{/Users/jim.maloney/github-clones/fragmented-aperture-book/images/appendixB/bootstrap_evals_hc168.pdf}
\caption{Bootstrap-resampled cumulative evaluations to first-hit
Good Enough proxy across the 168-bit Phase~1 HC pool. Each curve
is one of the six reference bands; the bootstrap median is
1{,}000--2{,}000 evals on the easier bands and 3{,}000--5{,}000 on the
harder ones --- not the $10^5$--$10^6$ that a naive
dimension-count argument would predict.}
\label{fig:bootstrap_evals_hc168}
\end{figure}

\begin{figure}[h]
\centering
\includegraphics[width=\linewidth]{/Users/jim.maloney/github-clones/fragmented-aperture-book/images/appendixB/convergence_hc168.pdf}
\caption{Per-trial convergence trajectories for the six 168-bit
reference bands. Each trace is one random-restart HC trial; the
y-axis is fitness, the x-axis is unique-design FDTD evaluations. The
trajectories share the same qualitative shape as 20-bit and 24-bit
HC: rapid early improvement, exponential slowdown into the local
basin, plateau at the band's reachable optimum.}
\label{fig:convergence_hc168}
\end{figure}

The finding is that 168-bit fragmented aperture optimization is
\emph{easier} than 24-bit optimization at the same physical
aperture, not harder. The reason is that more bits at a fixed
aperture means a finer pixelation of the same physical geometry,
not a more complex problem. Single-bit flips at 168~bits are
smaller perturbations than at 24~bits (a $\sim\!2.5$~cm patch vs
$\sim\!5$~cm at the apertures studied). Smaller perturbations
mean smoother landscape, more escape routes, and lower per-pass
fitness improvement variance. The curse of dimensionality, in
short, does not apply.

\subsection{The Curse of Dimensionality Does Not Apply}

The standard intuition about high-dimensional search --- that the
problem becomes harder as dimension grows --- assumes that each
additional dimension adds independent structure. For fragmented
apertures at fixed physical aperture, each additional bit refines
the resolution of an existing physical degree of freedom; the bit
encodes spatial fidelity, not new variation. The landscape's
intrinsic complexity is determined by the physics, not the bit
count. Increasing bit count gives the optimizer finer-grained moves
on the same underlying landscape.

This is the same high-dimensional phenomenon that makes stochastic
gradient descent practical in deep learning despite millions of
parameters: true local minima are exponentially rare in high
dimensions; most critical points are saddles with escape routes,
and a sequential greedy method (SGD or HC) traverses the saddle
structure without difficulty. The deep-learning analogy is
informal here, but the underlying mechanism is the same: each
additional dimension is more likely to be an escape route than a
trap, so dimension is a friend not an enemy.

\subsection{When the Classical-Method Comparison Matters in a Real Pipeline}

The practical pipeline at production scale is: surrogate forward
network for fitness prediction (Ch.~M), inverse network for
warm-start design generation (Ch.~M), then classical HC refinement
from the inverse-network seed to polish the design to a deployable
local optimum. The classical-method comparison of this appendix
matters for the third stage: HC is the right choice for refinement
because it is sequential, deterministic, cheap, and well-understood.
The two-criteria framework matters for evaluating whether the HC
refinement has converged to a deployable design.

At 168~bits and beyond, classical methods will rarely be used as
the primary search; the pipeline above is the production path. But
when the inverse network proposes a seed that misses the basin (the
mid-band failure mode of \cite{Maloney2026Paper1}), the classical
HC fallback is what catches the failure. Understanding the
fallback's behavior --- including when to escalate to SS-GA or
memetic on the surrogate (Ch.~M) --- is what this appendix
provides.

\section{What This Means for the Practicing Antenna Designer}
\label{app:practitioner}

\subsection{If You Have $<\!30$ Bits and One Band: Use HC, No Tuning Needed}

HC has no recipe to tune (the only knobs are random seed and
neighborhood size, and the default single-bit-flip neighborhood is
the right choice). At $<\!30$~bits HC will reach a local optimum
within $\sim\!60$--$300$ evaluations and the bootstrap-median
evals-to-Good-Enough is competitive with the best-tuned GA. There
is no reason to spend hyperparameter sweep budget unless the band
is in the frustrated regime (1.2--2.0~GHz at the apertures studied),
in which case see the next subsection.

\subsection{If You Have $<\!30$ Bits and Frustrated Mid-Bands: Use Memetic}

The frustrated mid-bands have target-basin coverage below 3\% and
memetic is the only method that reliably finds the basin. The
$\sim\!30\times$ per-trial cost over HC is worth it because the
trial \emph{success rate} is 60--95\% rather than 1--5\%, and the
cost-per-recovered-basin is essentially identical to HC's. Memetic
also has no tuning to speak of: the SS-GA core uses the
cached-best universal recipe of Sec.~\ref{app:recipes}, and the HC
polish on every offspring has no parameters.

\subsection{If You're Scaling to 100+ Bits Without Ground Truth: Use Surrogate Warm-Start}

Above $\sim\!100$~bits the right pipeline is the surrogate-driven
inverse-warm-start of Ch.~M, not classical methods at all. The
classical methods discussed in this appendix are then the HC
refinement step at the end of the pipeline (and the SS-GA or
memetic fallback on the surrogate when the warm-start seed misses
the basin). The cost balance has shifted at this scale: training a
forward surrogate costs $\sim\!100{,}000$ FDTD evaluations
(amortized across an entire program of design platforms); a single
warm-start design then costs $\sim\!60$ FDTD evaluations for the
HC polish. Classical-only methods at this scale would cost
$\sim\!5{,}000$ to $\sim\!1{,}000{,}000$ evaluations per design
(per Sec.~\ref{app:results} and \cite{Maloney2026Paper3}).

\subsection{Reporting Checklist}

For published optimizer comparisons in this literature, the
checklist that this appendix's methodology suggests is the
following. Without each item, comparisons are unreliable.
\begin{enumerate}
\item Cache accounting: state whether a seen-set is in use and what
      the cache hit rate is for each method. Without this, reported
      eval counts are ambiguous up to a factor of 2.
\item Bootstrap CI: report the bootstrap median and 5/95 percentile
      range, not just a point estimate. A single seed's
      evals-to-first-hit is not the metric.
\item Both criteria: report Good Enough and strict numbers
      separately. They answer different questions and a single
      headline number is not enough.
\item Win-count summary: across some chosen reference-band set,
      report which method wins each band. A single best-method
      claim is not defensible across the natural band variation.
\item Restart-budget extension: when the bootstrap-median is
      undefined (most trials do not reach the target), state the
      cumulative-cold-multistart cost instead. Reporting only the
      successful trials is sample-selection bias.
\end{enumerate}

The 20-bit and 24-bit databases are publicly available and
straightforward to use as a methodology check on any new classical
optimizer claim in the fragmented-aperture literature; we encourage
their use as a standard benchmark.

% ─────────────────────────────────────────────────────────────────────
\begin{thebibliography}{99}

\bibitem{Maloney1999URSI}
J.~G. Maloney, P.~H. Harms, M.~P. Kesler, T.~L. Fountain, and G.~S. Smith,
``Novel, planar antennas designed using the genetic algorithm,'' in
\emph{USNC/URSI National Radio Science Meeting}, Orlando, FL, 1999, p.~237.

\bibitem{Maloney2013Genetic24bit}
J.~G. Maloney, R.~T. Lee, and D.~W. Landgren,
``Genetic algorithms for fragmented aperture antennas: a complete
evaluation of a 24-bit design,'' in \emph{2013 USNC-URSI Radio Science
Meeting}, 2013, pp.~115.

\bibitem{Haupt1995Intro}
R.~L. Haupt, ``An introduction to genetic algorithms for
electromagnetics,'' \emph{IEEE Antennas and Propagation Magazine},
vol.~37, no.~2, pp.~7--15, 1995.

\bibitem{Haupt2007GA}
R.~L. Haupt and D.~H. Werner, \emph{Genetic Algorithms in
Electromagnetics}. Wiley-IEEE Press, 2007.

\bibitem{Goldberg1989GA}
D.~E. Goldberg, \emph{Genetic Algorithms in Search, Optimization, and
Machine Learning}. Addison-Wesley, 1989.

\bibitem{CarrollGA}
D.~L. Carroll, ``FORTRAN Genetic Algorithm Driver,'' 2001,
\url{http://www.cuaerospace.com/Technology/GeneticAlgorithm/}.

\bibitem{Moscato1989Memetic}
P.~Moscato, ``On evolution, search, optimization, genetic algorithms and
martial arts: towards memetic algorithms,'' Caltech Concurrent Computation
Program, Report 826, 1989.

\bibitem{DEAP}
F.-A. Fortin, F.-M. De~Rainville, M.-A. Gardner, M.~Parizeau, and C.~Gagn\'e,
``DEAP: Evolutionary algorithms made easy,'' \emph{Journal of Machine
Learning Research}, vol.~13, pp.~2171--2175, 2012.

\bibitem{Yee1966FDTD}
K.~S. Yee, ``Numerical solution of initial boundary value problems
involving Maxwell's equations in isotropic media,'' \emph{IEEE Trans.
Antennas Propag.}, vol.~14, no.~3, pp.~302--307, 1966.

\bibitem{Taflove2005CompEM}
A.~Taflove and S.~C. Hagness, \emph{Computational Electrodynamics: The
Finite-Difference Time-Domain Method}, 3rd~ed. Artech House, 2005.

\bibitem{Maloney2026Paper1}
J.~G. Maloney and S.~F. Mercer, ``AI-Accelerated Design of Fragmented
Aperture Antennas: Optimizer Benchmarks and Surrogate-Assisted Inverse
Design from an Exhaustive Million-Design Ground-Truth Database,''
\emph{Progress in Electromagnetics Research}, in submission, 2026.

\bibitem{Maloney2026Paper2}
J.~G. Maloney and S.~F. Mercer, ``AI-Accelerated Design of Fragmented
Aperture Antennas: 24-Bit Exhaustive Study,'' \emph{Progress in
Electromagnetics Research}, in submission, 2026.

\bibitem{Maloney2026Paper3}
J.~G. Maloney, ``168-bit fragmented aperture antenna design with
AI-accelerated optimization,'' in preparation, 2026.

\end{thebibliography}
